\section{K-018: the proposed maximal two-power determinant order is false}
\label{sec:amplifying-gadget}

\begin{theorem}[K-018]
For every integer $n\ge2$ there is a polystable tensor in
$(\mathbb C^n)^{\otimes3}$ whose diagonal stabilizer has determinant-product
image
\[
       \mu_{D_n},\qquad
       D_n=2^{\,n-1+\lfloor(n-2)/7\rfloor}.
                                                               \tag{K-018a}
\]
Its diagonal stabilizer is the product of the scalar two-torus and
a cyclic group of order $D_n$, and its determinant character is
faithful on that cyclic factor.  In particular, at $n=9$ its
determinant image has order512, greater than $2^{n-1}=256$.
Thus $2^{n-1}$ is not a universal upper bound, even for polystable
tensors with free, rationally independent supports.
\end{theorem}

\begin{proof}
We first give an exact eight-dimensional gadget.  Its distinguished
boundary triple is $(4,0,1)$.  In the order displayed below let $E$
be the following22 edges, and give edge $e_j$ mass
$(p_j+u_jr)/256$:
\[
\begin{gathered}
E=(262,540,056,336,755,517,622,420,303,114,275,472,\\
\hfill746,524,065,033,607,121,347,663,611,770),\\
p=(120,112,131,126,125,96,148,12,-15,208,136,-12,\\
\hfill-1,48,-5,130,15,48,145,141,-48,132),\\
u=(72,16,-43,-14,43,-96,-244,84,87,-80,-72,172,\\
\hfill57,80,29,14,169,80,-73,-101,176,-100).
\end{gathered}
\]
The one-coordinate margins are $r$ at the three boundary coordinates
and1 at every other coordinate.  This follows from the integer
identities $Ap=256({\bf1}-b)$ and $Au=256b$, where $A$ is the
24-by22 incidence matrix and $b$ has ones at the boundary triple.
The support is free.  Its incidence minor obtained by deleting the
boundary rows in the second and third modes has determinant256.
All these finite integer identities are verified by the retained
independent certificate script.

Every mass is strictly positive on
\[
                         I=(3/11,37/61).                      \tag{K-018b}
\]
Indeed each endpoint value is nonnegative, and no affine mass
vanishes at both endpoints.  Equivalently,
$11p+3u\ge0$ and $61p+37u\ge0$ componentwise, with no simultaneous
zero entries.  The distinguished output edge056 has mass
\[
             F(r)=\frac{131-43r}{256},\qquad
             F(I)=(25/61,41/88)\subset I.                    \tag{K-018c}
\]
In particular $1/2\in I$.

Here is the precise gluing operation.  Let an old free support have
positive margin-one masses and an edge $(i,j,k)$ of mass $r\in I$.
Delete that edge.  Introduce seven new coordinates in each mode,
identify the gadget boundary coordinates with $i,j,k$, and insert
the22 gadget edges with their displayed masses.  All old remaining
masses are unchanged.  The boundary margins replace precisely the
deleted mass $r$, and the new coordinates have margin1.  Thus the
new support is balanced and positive.  It is free because two old
edges already obey that condition, two gadget edges obey it, and
an old and a gadget edge could share two coordinates only at two
boundary coordinates; the unique old edge with that pair was
deleted.

We also need exact lattice control, rather than a numerical
denominator observation.  Suppose the old support in dimension
$a$ has $3a-2$ independent columns and maximal incidence minor of
absolute value $d$.  For a gadget in dimension $b$ with $3b-2$
independent columns and analogous index $d'$, this gluing has
$3(a+b-1)-2$ independent columns and index $dd'$.  To verify this,
delete the shared boundary row in each of the second and third
modes when taking the global maximal minor.  Delete the selected
old edge column from the old minor.  That old column had a single1
in the remaining boundary row of the first mode.  Order the old
rows except that row first, and all gadget rows last.  The global
matrix is block triangular: its upper right block is zero, its
upper left block is the old minor with that boundary row and edge
column removed, and its lower right block is the gadget minor.
Expansion in the old single1 column shows that its determinant has
absolute value $dd'$.  Complementary-minor duality from K-009
shows this is the torsion index of the full incidence lattice.

Start with the four edges000,011,101,110 in dimension2, each of
mass $1/2$.  Their incidence index is2.  Glue the gadget at000,
then repeatedly at its distinguished output edge056 (with the
indices transported by the gluing).  At the $k$th stage, dimension
and incidence index are
\[
                     n_k=2+7k,\qquad D_k=2\cdot256^k.
\]
Let $r_k$ be the mass of the selected edge.  Then
$r_0=1/2$ and $r_{k+1}=F(r_k)$, so every gluing is positive by
(K-018c).  Inductively $r_k$ has reduced denominator $D_k$ and
odd numerator: if $r_k=a/D_k$ with $a$ odd, then
\[
           r_{k+1}=\frac{131D_k-43a}{256D_k}
\]
has odd numerator.  All other masses have denominator dividing
$D_k$, by the affine gadget formula.  Consequently the exact
common denominator of the complete mass vector is $D_k$.

For independent incidence columns, the unique solution of
$Aq={\bf1}$ has denominator equal to the order of the determinant
character in $\mathbb Z^{3n}/\operatorname{im}A$:
$D{\bf1}$ lies in the incidence lattice if and only if every entry
of $Dq$ is integral.  Thus this character has order $D_k$.
The torsion subgroup itself also has order $D_k$ by the determinant
calculation.  It is therefore cyclic and generated by that
character, proving the claimed diagonal-stabilizer description.

To obtain every intermediate dimension, write
$n=2+7k+\ell$ with $0\le\ell\le6$.  Apply the ordinary edge-doubling
operation K-006 exactly $\ell$ times, each time selecting an edge
whose mass numerator is odd at the current exact common denominator.
That operation preserves positivity and freeness, multiplies the
incidence index and the selected denominator by2, and adds one
dimension.  The resulting order is
$2^{1+8k+\ell}=2^{n-1+k}$, as in (K-018a).

Finally set the tensor coefficient at each support edge equal to
the positive square root of its mass.  The one-coordinate reduced
Hermitian matrices have diagonal entries1.  Their off-diagonal
entries vanish by freeness.  Hence all three reduced matrices are
the identity and the tensor is critical.  Kempf--Ness, in the form
used and rechecked in P-F60, gives a closed $SL_n^3$ orbit.
This proves polystability and completes the theorem.
\end{proof}

This counterexample changes the proposed all-dimensional route:
the order $2^{n-1}$ is not the maximal two-power order.  K-018 by
itself gives no pole of the global invariant Hilbert series at its
new orders; classification and cancellation remain necessary for
that inference.  It does not refute a polynomial bound for
$p(k_n)$ or $s(k_n)$ by itself.

This is an explicit construction with no historical priority claim.
The primary-source scope audit is retained in
\texttt{reviews/axial\_denominator\_primary\_scope\_audit.md};
the precise denominator construction in the 2008 EJOR paper remains
an unread comparison gate.
